%%%%%%%%%%  scan idx 75  |  folio 62  %%%%%%%%%%
$\Atil$ (où $A$ est un anneau gradué), compatible avec les filtrations, et tel que l'homomorphisme correspondant $G(K(X))\longrightarrow G(A)$ soit injectif.}

Pour le prouver le lemme, nous utiliserons la théorie de Chow des classes de Chern~; si l'on pouvait démontrer le lemme, ou le théorème 2.14 directement, on en déduirait une théorie purement «~faisceautique~» des classes de Chern (cf.\ \S 5).

Nous considérerons l'homomorphisme $\Ctil$ de $\underline K(X)$ dans $\Atil(X)$ et nous poserons $A=A(X)$~; il faut prouver que l'homomorphisme $\psi=G(\Ctil)$ de $G(\underline K(X))$ dans $G(A(X))\simeq A(X)$ est injectif. Le fait que $\Ctil$ soit compatible avec les filtrations est vrai pour toute variété quasi-projective non singulière et résulte facilement du corollaire à la proposition 2.3. D'autre part comme $X$ est un produit de grasmaniennes $C_j$, on a sur $X$ des fibrés vectoriels $E_j$, images réciproques des fibrés canoniques sur les facteurs. Nous utiliserons les faits suivants (connus)~:
\begin{enumerate}[label=\alph*)]
\item Sur $X$, l'équivalence linéaire des cycles équivaut à l'équivalence numérique, et $A(X)$ est un groupe de rang fini~;
\item Les $C^i(E_j)$ engendrent l'anneau $A(X)$.
\end{enumerate}

De a) résulte que $\varphi:A(X)\longrightarrow G(\underline K(X))$ est injective (donc bijective) en vertu du lemme suivant~:

\textbf{Lemme 2.15.} \emph{Soient $X^n$ une variété projective non singulière, $Z^{n-p}$ un cycle sur $X$ tel que $\varphi(Z)\in G^p(\underline K(X))$ soit nul~; alors $Z^{n-p}$ est numériquement équivalent à zéro.}

Pour démontrer ce lemme, il suffit d'utiliser la fonction additive $\chi(F)=\sum_{i\geq0}(-1)^i\dim H^i(X,F)$ et le fait que cette fonction est égale à 
